% Archivo: sections/1_intro.tex
\section{Introducción}

En esta práctica se analiza la fenomenología del caos determinista a través de tres sistemas paradigmáticos: el péndulo no lineal forzado y amortiguado, el atractor de Lorenz y el péndulo doble.

El objetivo principal de este trabajo es caracterizar las transiciones hacia el caos y la topología del espacio de fases asociado a cada sistema. Para ello, se emplea integración numérica para resolver las ecuaciones diferenciales acopladas, permitiendo la extracción de métricas fundamentales como los diagramas de bifurcación, las secciones de Poincaré, el espectro de potencias y los exponentes de Lyapunov. Estas herramientas computacionales permiten distinguir rigurosamente entre regímenes periódicos, cuasiperiódicos y caóticos, estableciendo los límites de predictibilidad del sistema.